\begin{table}[tp]\centering\small
\caption{Four Calibrations of the Index Weights ($r^*$-Gap Yield and FCI; 12 Lags)}\label{tab:weights4}
\resizebox{\ifdim\width>\textwidth\textwidth\else\width\fi}{!}{%
\begin{tabular}{lcccccc}\toprule & Dollar relative & Window & Ordinary-day & \multicolumn{2}{c}{Yield and FCI, 24 m} & Shift in unemployment\\ \cmidrule(lr){5-6}
Weights & to equities & loading & loading & Unemployment & Price level & from adding the FCI, 24 m\\ \midrule
Market, futures factor (baseline) & 1.00 & 0.99 & 0.28 & 0.50$^{*}$ & $-$0.81$^{**}$ & 0.56 [0.10, 1.06]\\
Market, two-year yield & 2.06 & 0.99 & 0.65 & 0.39 & $-$0.66$^{**}$ & 0.46 [0.04, 0.87]\\
FRB/US & 2.51 & 0.99 & 0.37 & 0.56$^{*}$ & $-$0.71$^{**}$ & 0.63 [0.14, 1.14]\\
Goldman Sachs, published & 0.96 & 0.99 & 0.18 & 0.52$^{*}$ & $-$0.84$^{**}$ & 0.59 [0.12, 1.11]\\
\bottomrule\end{tabular}}
\begin{flushleft}\footnotesize\textit{Notes:} Per +1 pp on the two-year yield in the announcement window, 1996--2026. Each row builds the index from the same three $r^*$-gap series with the weights of one calibration in the companion paper \citep{fcinote}: market weights estimated with the futures factor or the two-year yield as the policy signal, FRB/US funds-rate equivalents, and the published Goldman Sachs weights; all are scaled to the same policy-news direction. Dollar relative to equities: the weight on a 3.9 percent dollar appreciation relative to that on a 5 percent fall in stock valuations. Ordinary-day loading: the FCI movement that accompanies a 1 pp rise in the yield on days without an announcement. Shift: the response with the FCI minus the response with the yield alone (same draws), with its 90 percent interval. $^{**}$ and $^{*}$: the 95 and 90 percent intervals exclude zero.\end{flushleft}
\end{table}
